% ----------------------------------------------------------
% Capítulo 3 — Metodologia
% ----------------------------------------------------------
\chapter{Metodologia}
\label{cap:metodologia}

Este capítulo descreve os procedimentos empíricos adotados para investigar o
impacto do método de amostragem --- \emph{dollar bars} versus \emph{time
bars} --- sobre a qualidade preditiva de modelos ARMA-GARCH para a variância
condicional de retornos do par BTC/USDT. A exposição segue a ordem
cronológica da \emph{pipeline} de modelagem: coleta e preparação dos dados
(\autoref{sec:met_dados}), análise exploratória das séries
(\autoref{sec:met_exploratoria}), construção dos retornos
(\autoref{sec:met_retornos}), divisão treino-teste
(\autoref{sec:met_split}), seleção de parâmetros via CPCV
(\autoref{sec:met_cpcv}), estimação do modelo final
(\autoref{sec:met_estimacao}), agregação das previsões para nível diário
(\autoref{sec:met_agregacao}), avaliação por métricas e testes estatísticos
(\autoref{sec:met_metricas}) e implementação computacional
(\autoref{sec:met_implementacao}).

% ==========================================================
\section{Dados e Período Amostral}
\label{sec:met_dados}
% ==========================================================

O estudo utiliza dados de negociação do par BTC/USDT na \emph{exchange}
Binance, abrangendo o período de 1\textordmasculine\ de janeiro de 2024 a
31 de dezembro de 2025 (731 dias corridos). A partir dos dados brutos de
alta frequência, foram construídas duas séries de barras com resolução
efetiva equivalente a 5 minutos:

\begin{itemize}
    \item \textbf{Dollar Bars (5 min)}: 209.630 barras, amostradas sempre
          que o volume financeiro acumulado atinge um limiar
          $\bar{D}$ calibrado para produzir, em média, o mesmo número de
          barras por dia que as \emph{time bars} de 5 minutos;
    \item \textbf{Time Bars (5 min)}: 210.522 barras, amostradas em
          intervalos fixos de 5 minutos.
\end{itemize}

Adicionalmente, a volatilidade realizada diária, estimada pelo método
TSRV (\autoref{subsec:tsrv}), serve como \emph{benchmark} para avaliação
das previsões. O \emph{benchmark} cobre os mesmos 731 dias do período
amostral.

% ----------------------------------------------------------
\subsection{Coleta e Pré-processamento dos Dados de Alta Frequência}
\label{subsec:met_coleta}
% ----------------------------------------------------------

Os dados brutos utilizados neste trabalho consistem nos registros de negociação agregada (\emph{aggTrades}) do par BTC/USDT disponibilizados pela Binance. Cada observação contém o preço de execução da negociação, o volume transacionado e a direção do \emph{taker} (campo \texttt{IsBuyerMaker}). O volume total do período amostral supera 1 bilhão de registros, distribuídos em arquivos no formato Apache Parquet com tamanho agregado de aproximadamente 28\,GB.

Dado o volume de dados incompatível com o carregamento integral em memória principal (\emph{RAM}), adotou-se uma arquitetura de processamento por \emph{streaming}, descrita em detalhe a seguir.

\subsubsection{Arquitetura de Processamento por \emph{Streaming} via \emph{Row Groups}}
\label{subsubsec:streaming_rg}

O formato Parquet organiza os dados em unidades lógicas denominadas \emph{row groups} --- blocos contíguos de linhas armazenados de forma colunar. A biblioteca \texttt{pyarrow} permite a leitura seletiva de um \emph{row group} por vez, sem que o restante do arquivo seja desserializado. Essa característica foi explorada para implementar um laço de processamento iterativo: a cada iteração, um único \emph{row group} é carregado em memória, as métricas de microestrutura são calculadas e acumuladas em estruturas de dados de tamanho fixo, e o bloco é então descartado antes que o próximo seja lido.

O resultado prático é um consumo de memória aproximadamente constante ao longo de toda a execução --- empiricamente observado em torno de 500\,MB --- independentemente do tamanho total do arquivo. Tal abordagem é condição necessária para a replicabilidade do experimento em ambientes com recursos computacionais convencionais (ausência de infraestrutura de \emph{cluster}), garantindo ao mesmo tempo que a íntegra dos dados seja processada sem amostragem prévia.

\subsubsection{Extração do Preço Eficiente por Filtro IIR}
\label{subsubsec:filtro_iir}

No arcabouço de microestrutura, o preço observado $P_t$ é decomposto em duas componentes:
%
\begin{equation}
    \label{eq:price_decomp}
    P_t = P_t^* + \varepsilon_t,
\end{equation}
%
onde $P_t^*$ denota o \emph{preço eficiente} --- o valor fundamental latente do ativo, livre de ruído transacional --- e $\varepsilon_t$ representa o ruído de microestrutura, gerado pela discrição do \emph{tick}, pelo \emph{spread} \emph{bid-ask} e por outros custos de transação.

Para extrair $P_t^*$, aplicou-se um filtro de Média Móvel Exponencial (EMA) de primeira ordem, equivalente a um filtro de Resposta ao Impulso Infinita (IIR) passa-baixas com função de transferência:
%
\begin{equation}
    \label{eq:iir}
    H(z) = \frac{\alpha}{1 - (1-\alpha)\,z^{-1}},
\end{equation}
%
onde $\alpha \in (0,1)$ é o coeficiente de suavização. A equação de recorrência correspondente é $P_t^* = \alpha P_t + (1-\alpha)P_{t-1}^*$, que pode ser implementada de forma causal e eficiente via \texttt{scipy.signal.lfilter}. A denominação ``IIR'' sublinha que, ao contrário dos filtros de Resposta ao Impulso Finita (FIR), cada saída depende de toda a história passada da série, tornando a manutenção do \emph{estado} do filtro entre os \emph{row groups} processados indispensável para garantir a continuidade do sinal: ao final de cada bloco, o vetor de estado interno do filtro é preservado e injetado como condição inicial do bloco subsequente.

O componente de ruído é então estimado como o resíduo $\hat{\varepsilon}_t = P_t - \hat{P}_t^*$, e suas propriedades estatísticas --- média, variância, percentis e perfil intradiário --- são acumuladas nas estruturas \emph{online} descritas nas subseções anteriores.

\subsubsection{Aceleração por Compilação JIT via Numba}
\label{subsubsec:numba}

Os \emph{kernels} de cálculo que envolvem laços internos sobre janelas móveis --- em particular, o estimador de Hasbrouck e a covariância serial do modelo de Roll --- foram decorados com a diretiva \texttt{@numba.njit} (modo \emph{nopython}), que instrui o compilador Numba a traduzir o código Python para código de máquina LLVM sem intermediação do interpretador CPython. O resultado é desempenho equivalente ao de C para laços numéricos intensivos, com o benefício de manter a legibilidade e a verificabilidade do código em Python. O ganho de velocidade observado em relação ao NumPy puro para janelas de comprimento $\sim$$10^4$ foi superior a uma ordem de grandeza, tornando o processamento da série completa viável em tempo compatível com o ciclo de pesquisa.

% ----------------------------------------------------------
\subsection{Estimação da Volatilidade Realizada}
\label{subsec:met_rv_estimacao}
% ----------------------------------------------------------

A Volatilidade Realizada (RV) é uma medida não paramétrica da variância integrada de um processo de preços, obtida pela soma dos quadrados dos log-retornos observados em alta frequência \cite{andersen2003}. Para um ativo com preços $P_0, P_1, \ldots, P_n$ amostrados a intervalos regulares, o estimador clássico é definido como:
%
\begin{equation}
    \label{eq:rv_classica}
    \widehat{RV}_n = \sum_{i=1}^{n} r_i^2,
    \qquad r_i = \ln\!\left(\frac{P_i}{P_{i-1}}\right).
\end{equation}
%
Sob o modelo de difusão sem ruído, $\widehat{RV}_n$ converge em probabilidade para a variância integrada $\int_0^1 \sigma_s^2\,\mathrm{d}s$ quando $n \to \infty$. No entanto, dados de altíssima frequência como \emph{aggTrades} da Binance são contaminados por ruído de microestrutura de mercado --- originado pelo \emph{bid-ask bounce}, pela discrição do \emph{tick} e por custos de transação ---, o que induz viés crescente em $\widehat{RV}_n$ à medida que a frequência de amostragem aumenta. Esse fenômeno é denominado o \emph{bias-variance trade-off} da volatilidade realizada.

% ----------------------------------------------------------
\subsubsection{Two-Scale Realized Volatility (TSRV)}
\label{subsec:tsrv}
% ----------------------------------------------------------

Para contornar o viés imposto pelo ruído de microestrutura, adotou-se o estimador \emph{Two-Scale Realized Volatility} (TSRV), proposto por \citeonline{zhang2005}. A ideia central é combinar duas escalas de observação: uma escala ``grossa'' (sub-grids), menos suscetível ao ruído, e a escala completa, que serve de base para a correção do viés.

Sejam $n$ log-retornos calculados sobre preços amostrados a uma frequência regular. O estimador TSRV é construído em três etapas:

\begin{enumerate}
    \item \textbf{RV na escala completa.} Calcula-se a RV clássica sobre todos os $n$ retornos:
    %
    \begin{equation}
        \label{eq:rv_all}
        \widehat{RV}^{(\text{all})} = \sum_{i=1}^{n} r_i^2.
    \end{equation}

    \item \textbf{Divisão em $K$ sub-grids e RV média.} Os índices $\{1, \ldots, n\}$ são particionados em $K$ sub-grids, onde o sub-grid $k$ contém os índices $\{k, k+K, k+2K, \ldots\}$. Para cada sub-grid $k$, calcula-se a RV parcial $\widehat{RV}^{(k)}$ a partir dos retornos entre observações consecutivas dentro do sub-grid. A média das $K$ RVs parciais fornece:
    %
    \begin{equation}
        \label{eq:rv_subgrid}
        \overline{RV}^{(K)} = \frac{1}{K}\sum_{k=1}^{K} \widehat{RV}^{(k)}.
    \end{equation}

    \item \textbf{Correção de viés.} O estimador TSRV corrige o viés de ruído subtraindo uma fração da RV na escala completa:
    %
    \begin{equation}
        \label{eq:tsrv}
        \widehat{TSRV} = \overline{RV}^{(K)} - \frac{\bar{n}}{n}\,\widehat{RV}^{(\text{all})},
    \end{equation}
    %
    onde $\bar{n} = n/K$ é o tamanho médio de cada sub-grid. A intuição é que o termo $(\bar{n}/n)\,\widehat{RV}^{(\text{all})}$ estima a contribuição do ruído à RV na escala completa, que é proporcional ao número de observações.
\end{enumerate}

\textbf{Determinação do $K$ ótimo.} \citeonline{zhang2005} demonstram que o $K$ que minimiza o erro quadrático médio assintótico do estimador é:
%
\begin{equation}
    \label{eq:k_otimo}
    K^* = \left\lfloor n^{2/3} \right\rceil,
\end{equation}
%
onde $\lfloor\cdot\rceil$ denota o arredondamento para o inteiro mais próximo. Na implementação, $K$ é limitado superiormente por $\lfloor n/2 \rfloor$ para garantir que cada sub-grid contenha ao menos duas observações. O estimador final é truncado em zero ($\max(\widehat{TSRV},\,0)$) para garantir não-negatividade em amostras finitas.

% ----------------------------------------------------------
\subsubsection{Amostragem em Tempo Calendário}
\label{subsubsec:amostragem_calendario}
% ----------------------------------------------------------

Os \emph{ticks} brutos chegam em tempo de transação irregular, tornando necessária uma reestruturação temporal antes do cômputo da TSRV. Adotou-se a amostragem em tempo calendário (\emph{calendar-time sampling}): para cada frequência $\Delta$ configurável (ex.: 30s, 1 min, 2 min), o preço correspondente a cada intervalo $[t, t+\Delta)$ é definido como o preço da \emph{última} transação executada naquele intervalo (\texttt{resample($\Delta$).last()}). Intervalos sem negociação são descartados (\texttt{dropna()}). Esse procedimento produz uma grade temporal regular de preços $P_1, P_2, \ldots, P_n$ sobre a qual os log-retornos são calculados via diferenciação logarítmica: $r_i = \ln(P_i) - \ln(P_{i-1})$.

% ----------------------------------------------------------
\subsubsection{Processamento Diário Eficiente (\texttt{TickDataProcessor})}
\label{subsubsec:tick_processor}
% ----------------------------------------------------------

Dado o volume de dados (superior a 1 bilhão de registros), o processamento é realizado dia a dia. A classe \texttt{TickDataProcessor} encapsula a leitura seletiva de dados do arquivo Parquet via filtros aplicados diretamente no disco (\texttt{pyarrow.parquet.read\_parquet} com predicado de intervalo de datas), carregando em memória apenas as colunas \texttt{DateTime}, \texttt{Price} e \texttt{Qty} referentes a um único dia de negociação. Para cada dia, a função \texttt{compute\_daily\_rv} realiza: (i) a reestruturação temporal por reamostragem; (ii) a instanciação da classe \texttt{TwoScaleRV} sobre o vetor de preços resultante; e (iii) o cômputo do estimador com coleta de estatísticas auxiliares --- número de \emph{ticks}, número de observações amostradas, volume financeiro e preço médio.

% ----------------------------------------------------------
\subsubsection{Benchmark de Consenso Multi-Frequência}
\label{subsubsec:benchmark_consenso}
% ----------------------------------------------------------

A escolha da frequência de reamostragem $\Delta$ introduz um grau de arbitrariedade nas estimativas de TSRV. Para mitigar esse efeito, construiu-se um \emph{benchmark} de consenso: a \emph{pipeline} \texttt{compute\_rv\_series} é executada independentemente para três frequências --- 30 segundos, 1 minuto e 2 minutos --- gerando três séries diárias de RV, denotadas $\widehat{TSRV}_{30\text{s},d}$, $\widehat{TSRV}_{1\text{min},d}$ e $\widehat{TSRV}_{2\text{min},d}$. O benchmark de consenso para o dia $d$ é então definido como a média aritmética dessas três estimativas:
%
\begin{equation}
    \label{eq:rv_consenso}
    \widehat{RV}_{\text{consenso},d} = \frac{1}{3}\!\left(\widehat{TSRV}_{30\text{s},d} + \widehat{TSRV}_{1\text{min},d} + \widehat{TSRV}_{2\text{min},d}\right),
\end{equation}
%
com dias em que qualquer das três frequências produziu valor ausente (\texttt{NaN}) tratados de forma robusta pela média sobre os valores disponíveis. Essa abordagem reduz a sensibilidade do benchmark a escolhas arbitrárias de $\Delta$ e fornece uma medida mais estável da variância integrada diária.

A volatilidade anualizada correspondente é obtida por:
%
\begin{equation}
    \label{eq:vol_anualizada}
    \hat{\sigma}_{\text{anual},d} = \sqrt{\widehat{RV}_{\text{consenso},d} \times 252}.
\end{equation}

% ==========================================================
\section{Análise Exploratória das Séries}
\label{sec:met_exploratoria}
% ==========================================================

Antes de proceder à modelagem, realizou-se uma análise exploratória das
séries completas de retornos logarítmicos, com o objetivo de verificar as
premissas dos modelos ARMA-GARCH e identificar diferenças distribucional
e de dependência entre os dois métodos de amostragem.

% ----------------------------------------------------------
\subsection{Estatísticas Descritivas}
\label{subsec:met_descritivas}
% ----------------------------------------------------------

A \autoref{tab:descritivas} apresenta os quatro primeiros momentos
amostrais da distribuição dos retornos para ambos os tipos de barra. As
estatísticas foram calculadas sobre a série completa (209.630 observações
para \emph{dollar bars} e 210.522 para \emph{time bars}).

\begin{table}[H]
\centering
\caption{Estatísticas descritivas dos retornos logarítmicos.}
\label{tab:descritivas}
\begin{tabular}{lSS}
\toprule
{Estatística} & {\emph{Dollar Bars}} & {\emph{Time Bars}} \\
\midrule
Média ($1^{\circ}$ momento)        &  3.46e-6  &  3.44e-6  \\
Variância ($2^{\circ}$ momento)    &  2.49e-6  &  2.40e-6  \\
Assimetria ($3^{\circ}$ momento)   &  0.0547   & -1.0071   \\
Curtose de Fisher ($4^{\circ}$ momento) &  3.0879   &  75.3687  \\
\bottomrule
\end{tabular}
\fonte{Elaboração própria.}
\end{table}

Dois aspectos merecem destaque. Primeiro, ambas as séries apresentam
excesso de curtose positivo, indicando caudas mais pesadas do que a
distribuição normal --- um fato estilizado amplamente documentado em
retornos financeiros \cite{tsay2010}. Entretanto, a magnitude do excesso
de curtose é drasticamente diferente: enquanto os retornos de
\emph{dollar bars} exibem curtose de Fisher igual a 3,09, os retornos de
\emph{time bars} alcançam 75,37, evidenciando a presença de observações
extremas significativamente mais severas.

Segundo, a assimetria dos retornos de \emph{time bars} ($-1,01$) é
substancialmente mais negativa do que a dos \emph{dollar bars} ($0,05$,
praticamente simétrica). Isso sugere que a amostragem temporal fixa
captura eventos de cauda esquerda (quedas abruptas) de forma mais
concentrada, possivelmente por comprimir muitas transações em uma única
barra durante episódios de alta volatilidade.

Esses resultados corroboram a tese de
\citeonline{lopez2018advances}: \emph{dollar bars} produzem retornos
com propriedades distribucional mais favoráveis à modelagem estatística.

% ----------------------------------------------------------
\subsection{Testes de Normalidade, Estacionariedade e Heteroscedasticidade}
\label{subsec:met_testes}
% ----------------------------------------------------------

A \autoref{tab:testes} sintetiza os resultados dos testes aplicados.

\begin{table}[H]
\centering
\caption{Testes de normalidade, estacionariedade e
heteroscedasticidade sobre os retornos logarítmicos.}
\label{tab:testes}
\begin{tabular}{lSSSS}
\toprule
{Teste} &
\multicolumn{2}{c}{\emph{Dollar Bars}} &
\multicolumn{2}{c}{\emph{Time Bars}} \\
\cmidrule(lr){2-3} \cmidrule(lr){4-5}
& {Estat.} & {\emph{p}-valor} & {Estat.} & {\emph{p}-valor} \\
\midrule
Jarque-Bera          & 83384.65  & {$<0{,}001$} & 49860565.37 & {$<0{,}001$} \\
Kolmogorov-Smirnov   & 0.0123    & {$<0{,}001$} & 0.0934      & {$<0{,}001$} \\
\addlinespace
ADF                  & -102.46   & {$<0{,}001$} & -57.32      & {$<0{,}001$} \\
KPSS                 & 0.2053    & 0.10         & 0.2171      & 0.10         \\
\addlinespace
Ljung-Box ($r_t$, lag 10)   & 50.06    & {$<0{,}001$} & 155.03   & {$<0{,}001$} \\
Ljung-Box ($r_t^2$, lag 10) & 79528.20 & {$<0{,}001$} & 28237.22 & {$<0{,}001$} \\
ARCH de Engle (lag 10)      & 30420.71 & {$<0{,}001$} & 24064.78 & {$<0{,}001$} \\
\bottomrule
\end{tabular}
\nota{H$_0$ por teste: Jarque-Bera e KS — normalidade;
ADF — raiz unitária; KPSS — estacionariedade;
Ljung-Box e ARCH de Engle — ausência de autocorrelação/efeito ARCH.}
\fonte{Elaboração própria.}
\end{table}

Os resultados permitem extrair três conclusões:

\begin{enumerate}
    \item \textbf{Não-normalidade}: ambos os testes de normalidade
          (Jarque-Bera e Kolmogorov-Smirnov) rejeitam a hipótese nula a
          qualquer nível de significância convencional ($p < 0{,}001$), confirmando
          que os retornos não seguem distribuição gaussiana. Nota-se que a
          estatística de Jarque-Bera para \emph{time bars} é mais de 500
          vezes maior do que para \emph{dollar bars}, refletindo a curtose
          extrema da amostragem temporal;
    \item \textbf{Estacionariedade}: o teste ADF rejeita a hipótese nula
          de raiz unitária para ambas as séries ($p < 0{,}001$), enquanto o
          teste KPSS \emph{não} rejeita a hipótese nula de estacionariedade
          ($p = 0{,}10$). Os dois testes são consistentes e confirmam que as
          séries de retornos são estacionárias, não sendo necessária
          diferenciação ($d = 0$);
    \item \textbf{Efeitos ARCH}: tanto o teste de Ljung-Box sobre os
          retornos quadráticos quanto o teste ARCH de Engle rejeitam
          fortemente a hipótese nula de ausência de heteroscedasticidade
          condicional ($p < 0{,}001$), justificando a utilização de modelos
          da família GARCH para capturar a dinâmica da variância
          condicional.
\end{enumerate}

Adicionalmente, a \autoref{fig:normalidade_dollar} e a
\autoref{fig:normalidade_time} apresentam os histogramas com curva normal
ajustada e os QQ-plots dos retornos para cada tipo de barra, ilustrando
visualmente o afastamento da normalidade.

\begin{figure}[H]
    \centering
    \includegraphics[width=0.95\textwidth]{Dollar_Bars_5m_normalidade.png}
    \caption{Histograma com curva normal e QQ-Plot --- Dollar Bars (5 min).}
    \label{fig:normalidade_dollar}
\end{figure}

\begin{figure}[H]
    \centering
    \includegraphics[width=0.95\textwidth]{Time_Bars_5m_normalidade.png}
    \caption{Histograma com curva normal e QQ-Plot --- Time Bars (5 min).}
    \label{fig:normalidade_time}
\end{figure}

A \autoref{fig:acf_dollar_sq} e a \autoref{fig:acf_time_sq} apresentam as
funções de autocorrelação (FAC) e autocorrelação parcial (FACP) dos
retornos ao quadrado, evidenciando a persistência da volatilidade
condicional em ambas as séries.

\begin{figure}[H]
    \centering
    \includegraphics[width=0.95\textwidth]{Dollar_Bars_5m_acf_pacf_retornos2.png}
    \caption{ACF e PACF dos retornos ao quadrado --- Dollar Bars (5 min).}
    \label{fig:acf_dollar_sq}
\end{figure}

\begin{figure}[H]
    \centering
    \includegraphics[width=0.95\textwidth]{Time_Bars_5m_acf_pacf_retornos2.png}
    \caption{ACF e PACF dos retornos ao quadrado --- Time Bars (5 min).}
    \label{fig:acf_time_sq}
\end{figure}

% ==========================================================
\section{Construção dos Retornos Logarítmicos}
\label{sec:met_retornos}
% ==========================================================

Para cada série de barras, os retornos logarítmicos foram calculados como:
%
\begin{equation}
    \label{eq:met_log_return}
    r_t = \ln\left(\frac{C_t}{C_{t-1}}\right),
\end{equation}
%
onde $C_t$ denota o preço de fechamento da $t$-ésima barra. A primeira
observação de cada série é descartada ($r_1$ é indefinido), resultando em
209.629 retornos para \emph{dollar bars} e 210.521 para \emph{time bars}.

A opção por retornos logarítmicos em vez de retornos simples
($R_t = C_t/C_{t-1} - 1$) justifica-se por duas razões: (i)~retornos
logarítmicos são aditivos ao longo do tempo, facilitando a agregação
temporal de variâncias; e (ii)~para retornos de pequena magnitude, como é
o caso de barras de 5 minutos, $r_t \approx R_t$, de modo que a escolha é
praticamente irrelevante numericamente.

% ==========================================================
\section{Divisão Treino-Teste}
\label{sec:met_split}
% ==========================================================

A amostra total de cada série foi dividida em dois conjuntos por meio de um
\emph{split} temporal na proporção 80\%/20\%, preservando a ordem
cronológica dos dados. A \autoref{tab:split} resume as dimensões e períodos
de cada partição.

\begin{table}[H]
\centering
\caption{Divisão treino-teste por tipo de barra.}
\label{tab:split}
\begin{tabular}{llrr}
\toprule
{Série} & {Partição} & {Observações} & {Período} \\
\midrule
\multirow{2}{*}{\emph{Dollar Bars}}
    & Treino (80\%) & 167.704 & 01/jan/2024 -- 14/jul/2025 \\
    & Teste (20\%)  &  41.926 & 14/jul/2025 -- 31/dez/2025 \\
\addlinespace
\multirow{2}{*}{\emph{Time Bars}}
    & Treino (80\%) & 168.417 & 01/jan/2024 -- 07/ago/2025 \\
    & Teste (20\%)  &  42.105 & 07/ago/2025 -- 31/dez/2025 \\
\bottomrule
\end{tabular}
\fonte{Elaboração própria.}
\end{table}

Observa-se que as datas de corte diferem ligeiramente entre os dois tipos
de barra (14 de julho vs. 7 de agosto de 2025), uma vez que a proporção
80\%/20\% é definida sobre o número de barras, e não sobre o calendário.
Isso é esperado, pois as \emph{dollar bars} possuem espaçamento temporal
irregular.

É importante ressaltar que a divisão treino-teste é distinta do
procedimento de validação cruzada CPCV descrito na
\autoref{sec:met_cpcv}: o CPCV opera exclusivamente sobre o conjunto de
\textbf{treino} para selecionar os hiperparâmetros, enquanto o conjunto de
\textbf{teste} é reservado para a avaliação final do modelo selecionado.

% ==========================================================
\section{Seleção de Parâmetros via CPCV}
\label{sec:met_cpcv}
% ==========================================================

A seleção das ordens ótimas $(p, q)$ para a equação da média
(ARMA) e $(r, s)$ para a equação da variância (GARCH) foi realizada pelo
arcabouço de Validação Cruzada Combinatória Purgada (CPCV), descrito na
\autoref{sec:cpcv}. Este procedimento foi aplicado exclusivamente ao
conjunto de treino de cada série.

% ----------------------------------------------------------
\subsection{Configuração do CPCV}
\label{subsec:met_cpcv_config}
% ----------------------------------------------------------

O CPCV foi configurado com os seguintes parâmetros:

\begin{itemize}
    \item \textbf{Número de grupos}: $N = 10$ (o conjunto de treino é
          dividido em 10 blocos sequenciais);
    \item \textbf{Grupos de teste}: $k = 5$ (a cada partição, 5 grupos
          compõem o teste);
    \item \textbf{Número de folds}: $\binom{10}{5} = 252$ partições
          treino-teste distintas;
    \item \textbf{Embargo}: 240 barras ($\approx$ 20 horas em barras de
          5 min) após cada bloco de teste, mitigando vazamento por
          autocorrelação serial.
\end{itemize}

Em cada fold, as observações situadas na fronteira entre os blocos de
treino e teste são removidas (purgação), e as 240 barras subsequentes ao
final de cada bloco de teste são embargadas, conforme o procedimento
descrito na \autoref{subsec:cpcv_formulacao}. Em média, cada fold utiliza
aproximadamente 81.500 observações para treino e 84.000 para teste,
com cerca de 1.200 observações purgadas e 1.200 embargadas.

% ----------------------------------------------------------
\subsection{Grid de Busca}
\label{subsec:met_grid}
% ----------------------------------------------------------

O espaço de busca compreende o produto cartesiano:
%
\begin{equation}
    \label{eq:met_grid}
    \text{ARMA}(p, q) \times \text{GARCH}(r, s),
    \quad p, q \in \{0, 1, \ldots, 5\},\;
    r \in \{1, 2\},\;
    s \in \{1, 2\},
\end{equation}
%
totalizando $6 \times 6 \times 2 \times 2 = 144$ combinações de
parâmetros. Para cada uma das 252 partições CPCV, todas as 144
combinações foram estimadas por máxima verossimilhança (MV), resultando em
$252 \times 144 = 36.288$ ajustes MLE por tipo de barra e 72.576 ajustes
no total da \emph{pipeline}.

O critério de seleção em cada fold é o AIC
(\autoref{eq:aic}): a combinação com o menor AIC no
conjunto de validação do fold é registrada como ``vencedora''. Ao final dos
252 folds, o modelo selecionado é aquele cuja combinação de parâmetros
obteve a maior \emph{frequência} de vitórias --- ou seja, a moda da
distribuição de modelos selecionados ao longo dos folds.

% ----------------------------------------------------------
\subsection{Resultados da Seleção}
\label{subsec:met_cpcv_resultados}
% ----------------------------------------------------------

A \autoref{tab:cpcv_resultados} apresenta os modelos selecionados para
cada tipo de barra.

\begin{table}[H]
\centering
\caption{Resultados da seleção de parâmetros via CPCV.}
\label{tab:cpcv_resultados}
\begin{tabular}{llrl}
\toprule
{Série} & {Modelo selecionado} & {Frequência} & {AIC médio} \\
\midrule
\emph{Dollar Bars} & ARMA(0,0)-GARCH(1,2) & 233/252 (92,5\%) &
    $-835.475$ \\
\emph{Time Bars}   & ARMA(0,0)-GARCH(1,1) & 252/252 (100\%) &
    $-859.083$ \\
\bottomrule
\end{tabular}
\fonte{Elaboração própria.}
\end{table}

Dois resultados merecem destaque:

\begin{enumerate}
    \item \textbf{Ausência de componente ARMA}: para ambos os tipos de
          barra, a equação da média selecionada foi ARMA(0,0) --- ou seja,
          média constante sem termos autorregressivos ou de médias móveis.
          Isso é consistente com a hipótese de mercados eficientes na forma
          fraca e com a evidência da \autoref{tab:testes}, na qual o teste
          de Ljung-Box detecta autocorrelação nos retornos, mas de
          magnitude insuficiente para justificar termos ARMA (as
          informações preditivas concentram-se nos retornos ao quadrado,
          captados pelo GARCH);
    \item \textbf{Ordem GARCH distinta}: a amostragem por \emph{dollar
          bars} requer um modelo GARCH(1,2) --- com dois termos
          de variâncias defasadas --- enquanto as \emph{time bars} são
          adequadamente descritas pelo GARCH(1,1) padrão. A ordem
          superior para \emph{dollar bars} sugere uma dinâmica de
          persistência de volatilidade ligeiramente mais complexa nessa
          amostragem, possivelmente refletindo a irregularidade temporal
          das barras.
\end{enumerate}

A \autoref{fig:cpcv_aic} apresenta a evolução do AIC médio por fold para
ambos os modelos selecionados, evidenciando a estabilidade da seleção ao
longo das 252 partições.

\begin{figure}[H]
    \centering
    \includegraphics[width=0.95\textwidth]{comparacao_modelos/cpcv_01_aic_por_fold.png}
    \caption{AIC médio por fold CPCV --- Dollar Bars vs.\ Time Bars.}
    \label{fig:cpcv_aic}
\end{figure}

% ==========================================================
\section{Estimação do Modelo Final}
\label{sec:met_estimacao}
% ==========================================================

Após a seleção de parâmetros, o modelo final de cada série foi
reestimado utilizando \textbf{todo} o conjunto de treino (sem
particionamento CPCV), a fim de maximizar a precisão dos coeficientes
estimados. A \autoref{tab:modelo_final} apresenta os critérios de
informação do ajuste final.

\begin{table}[H]
\centering
\caption{Critérios de informação do modelo final estimado sobre o
conjunto de treino completo.}
\label{tab:modelo_final}
\begin{tabular}{llSS}
\toprule
{Série} & {Modelo} & {AIC} & {BIC} \\
\midrule
\emph{Dollar Bars} & GARCH(1,2) & -1693794.21 & -1693744.06 \\
\emph{Time Bars}   & GARCH(1,1) & -1742053.06 & -1742012.92 \\
\bottomrule
\end{tabular}
\fonte{Elaboração própria.}
\end{table}

A partir do modelo estimado, realizou-se a previsão \emph{one-step-ahead}
da variância condicional $\hat{\sigma}_t^2$ para cada observação do
conjunto de teste: 41.926 previsões para \emph{dollar bars} e 42.105 para
\emph{time bars}.

% ==========================================================
\section{Agregação para Nível Diário}
\label{sec:met_agregacao}
% ==========================================================

Como a volatilidade realizada (RV) de referência é calculada em base
diária, as previsões intradiárias de variância condicional precisam ser
agregadas ao nível diário para permitir comparação direta. A agregação
segue o princípio de aditividade da variância de retornos logarítmicos
independentes: a variância diária prevista é obtida pela soma das
variâncias condicionais intradiárias dentro de cada dia:
%
\begin{equation}
    \label{eq:met_agregacao}
    \hat{\sigma}^2_{\text{dia},d} =
    \sum_{t \in \mathcal{B}_d} \hat{\sigma}_t^2,
\end{equation}
%
onde $\mathcal{B}_d$ denota o conjunto de barras pertencentes ao dia $d$.

Após a agregação e o alinhamento com o \emph{benchmark} de volatilidade
realizada, obtiveram-se:

\begin{itemize}
    \item \textbf{Dollar Bars}: 171 dias com previsão (14/jul/2025 a
          31/dez/2025), todos alinhados com o RV \emph{benchmark};
    \item \textbf{Time Bars}: 147 dias com previsão (07/ago/2025 a
          31/dez/2025), todos alinhados com o RV \emph{benchmark}.
\end{itemize}

A diferença na quantidade de dias (171 vs.\ 147) decorre das datas
distintas de início do conjunto de teste para cada tipo de barra,
conforme discutido na \autoref{sec:met_split}. Para garantir comparabilidade
estrita entre os modelos, a avaliação final das métricas, os intervalos
de confiança e o teste de Diebold-Mariano foram calculados exclusivamente
sobre a \textbf{interseção} dos dois períodos: os \textbf{147 dias comuns}
compreendidos entre 07/ago/2025 e 31/dez/2025.

% ==========================================================
\section{Métricas de Avaliação e Testes Estatísticos}
\label{sec:met_metricas}
% ==========================================================

As previsões diárias de variância foram avaliadas contra o \emph{benchmark}
de volatilidade realizada (RV consenso, \autoref{subsec:tsrv}) utilizando
as métricas e testes descritos na \autoref{sec:metricas}. Esta seção
detalha a parametrização específica empregada.

% ----------------------------------------------------------
\subsection{Funções de Perda e Regressão Mincer-Zarnowitz}
\label{subsec:met_funcoes_perda}
% ----------------------------------------------------------

Foram calculadas três funções de perda --- MSE (\autoref{eq:mse}), MAE
(\autoref{eq:mae}) e QLIKE (\autoref{eq:qlike}) --- sobre os pares
$(\hat{\sigma}^2_{\text{dia},d},\, \xi_d)$ para cada tipo de barra, onde
$\xi_d$ denota a volatilidade realizada diária do \emph{benchmark}.

A regressão de Mincer-Zarnowitz (\autoref{eq:mz}) foi estimada por MQO,
fornecendo o coeficiente de determinação $R^2$, o intercepto $\alpha$ e a
inclinação $\beta$, com seus respectivos $p$-valores.

A \autoref{tab:metricas} apresenta os resultados comparativos.

\begin{table}[H]
\centering
\caption{Métricas de avaliação das previsões de variância vs.\ RV
\emph{benchmark}.}
\label{tab:metricas}
\begin{tabular}{lSS}
\toprule
{Métrica} & {\emph{Dollar Bars}} & {\emph{Time Bars}} \\
\midrule
MSE               & 1.05e-7  & 1.65e-7  \\
MAE               & 2.64e-4  & 3.59e-4  \\
QLIKE             & -6.993   & -6.776   \\
$R^2_{\text{MZ}}$ & 0.667    & 0.493    \\
$\hat{\alpha}$    & -9.20e-5 & -2.54e-4 \\
$\hat{\beta}$     & 0.868    & 1.032    \\
\bottomrule
\end{tabular}
\fonte{Elaboração própria.}
\end{table}

% ----------------------------------------------------------
\subsection{Teste de Diebold-Mariano}
\label{subsec:met_dm}
% ----------------------------------------------------------

Para comparar formalmente o poder preditivo dos dois modelos, aplicou-se o
teste de Diebold-Mariano (\autoref{subsec:diebold_mariano}) com a função
de perda MSE. A hipótese nula é de igual poder preditivo entre o modelo
ajustado sobre \emph{dollar bars} e o ajustado sobre \emph{time bars}.

O teste foi calculado sobre o subconjunto de dias em que ambos os modelos
possuem previsões, com correção HAC de Newey-West. Os resultados foram:

\begin{itemize}
    \item Estatística DM: $-4{,}806$;
    \item $p$-valor: $< 0{,}001$.
\end{itemize}

Ao nível de significância de 5\%, \textbf{rejeita-se} a hipótese nula de
igual poder preditivo ($p < 0{,}001$). A estatística negativa indica que
o modelo ajustado sobre \emph{dollar bars} produz perdas MSE
significativamente menores do que o ajustado sobre \emph{time bars},
evidenciando superioridade preditiva estatisticamente distinguível.

% ----------------------------------------------------------
\subsection{Intervalos de Confiança Bootstrap}
\label{subsec:met_bootstrap}
% ----------------------------------------------------------

Intervalos de confiança de 95\% para as métricas pontuais foram
construídos pelo método de \emph{bootstrap} percentil não-paramétrico
(\autoref{subsec:bootstrap_ic}) com $B = 2.000$ reamostras. A
\autoref{tab:bootstrap} apresenta os resultados.

\begin{table}[H]
\centering
\caption{Intervalos de confiança bootstrap (95\%) para as métricas de
avaliação.}
\label{tab:bootstrap}
\begin{tabular}{llScc}
\toprule
{Série} & {Métrica} & {Pontual} & {IC 95\%} & {SE} \\
\midrule
\multirow{3}{*}{\emph{Dollar Bars}}
    & MSE   & 1.050e-7 & {$[7{,}34\text{e-}8;\; 1{,}49\text{e-}7]$} & 1.93e-8 \\
    & MAE   & 2.637e-4 & {$[2{,}35\text{e-}4;\; 2{,}97\text{e-}4]$} & 1.59e-5 \\
    & QLIKE & -6.993   & {$[-7{,}137;\; -6{,}845]$}                  & 7.48e-2 \\
\addlinespace
\multirow{3}{*}{\emph{Time Bars}}
    & MSE   & 1.645e-7 & {$[1{,}35\text{e-}7;\; 2{,}00\text{e-}7]$} & 1.66e-8 \\
    & MAE   & 3.588e-4 & {$[3{,}29\text{e-}4;\; 3{,}90\text{e-}4]$} & 1.56e-5 \\
    & QLIKE & -6.776   & {$[-6{,}877;\; -6{,}668]$}                  & 5.38e-2 \\
\bottomrule
\end{tabular}
\nota{Valores calculados sobre os 147 dias comuns (07/ago/2025--31/dez/2025).
$B = 2.000$ reamostras; IC pelo método percentil.}
\fonte{Elaboração própria.}
\end{table}

Os intervalos de confiança para o MSE \textbf{não se sobrepõem}
($[7{,}34\text{e-}8;\; 1{,}49\text{e-}7]$ vs.\
$[1{,}35\text{e-}7;\; 2{,}00\text{e-}7]$), corroborando a rejeição da
hipótese nula pelo teste de Diebold-Mariano. Para a QLIKE, a ausência de
sobreposição também é verificada, reforçando a superioridade das
\emph{dollar bars} sob essa função de perda robusta \cite{patton2011}.

